El tubo de Franck-Hertz debe calentarse para evaporar el mercurio y asegurar 
que tenga una energía similar a los electrones que se desprenden del cátodo.

Este tubo está equipado con dos rejillas (G) y tres potenciales: un potencial 
de preaceleración ($U_1$), que acelera a los electrones hasta alcanzar una 
energía cercana a 4.9 eV, correspondiente a la energía de excitación del 
mercurio; un potencial de aceleración ($U_2$), que varía en forma de rampa 
desde 1 hasta 30 V; y un potencial de frenado ($U_3$), ajustado para mejorar la 
visualización de los picos. El esquema de este tubo se puede observar en la 
\cref{fig:Franck-Hertz-tube}

\begin{figure}[htbp!]                                                           
  \begin{center}                                                                
    \includegraphics[height=0.7\linewidth, page=2]{../Images/Franck-Hertz-tube.jpg}
    \caption{Tubo de Franck-Hertz.}                                             
    \label{fig:Franck-Hertz-tube}                                              
  \end{center}                                                                  
\end{figure}

La corriente registrada en el ánodo se midió mediante un osciloscopio 
analógico, encargado de graficar la relación entre el potencial y la corriente.
Este montaje se puede observar en la \cref{fig:experimental-setup}.

\begin{figure}[htbp!]
  \begin{center}
    \includegraphics[height=0.7\linewidth, page=2]{../Images/experimental-setup.jpg}
    \caption{Montaje experimental.}
    \label{fig:experimental-setup}
  \end{center}
\end{figure}
